% !TeX program = lualatex
% Compilation : lualatex choix-bibliotheque-pdf.tex
\documentclass[11pt,a4paper]{article}

\usepackage{fontspec}
\usepackage{polyglossia}
\setmainlanguage{french}
\setotherlanguage{english}
\setmainfont{TeX Gyre Pagella}
\setsansfont{TeX Gyre Heros}
\setmonofont{DejaVu Sans Mono}

\usepackage[a4paper,top=2.1cm,bottom=2.1cm,left=2.1cm,right=2.1cm]{geometry}
\usepackage{microtype}
\usepackage[table]{xcolor}
\usepackage{booktabs}
\usepackage{tabularx}
\usepackage{longtable}
\usepackage{array}
\usepackage{enumitem}
\usepackage{tcolorbox}
\tcbuselibrary{skins}
\usepackage{pdflscape}
\usepackage{fancyhdr}
\usepackage{lastpage}
\usepackage{titlesec}
\usepackage{hyperref}

\definecolor{bleuDAE}{HTML}{174A73}
\definecolor{bleuClairDAE}{HTML}{EAF3F8}
\definecolor{vertDAE}{HTML}{2E7D5B}
\definecolor{vertClairDAE}{HTML}{EAF6F0}
\definecolor{orangeDAE}{HTML}{B76419}
\definecolor{orangeClairDAE}{HTML}{FFF1E3}
\definecolor{rougeDAE}{HTML}{9B3155}
\definecolor{grisTexteDAE}{HTML}{40464D}
\definecolor{grisFondDAE}{HTML}{F4F6F7}

\hypersetup{
  colorlinks=true,
  linkcolor=bleuDAE,
  urlcolor=bleuDAE,
  pdftitle={Aide à la décision — Bibliothèque PDF},
  pdfauthor={Projet GDAETF2},
}

\setlength{\parindent}{0pt}
\setlength{\parskip}{0.45em}
\setlist[itemize]{leftmargin=*,itemsep=0.2em,topsep=0.2em}
\setlist[enumerate]{leftmargin=*,itemsep=0.25em,topsep=0.2em}
\titleformat{\section}{\Large\bfseries\color{bleuDAE}}{\thesection}{0.6em}{}
\titleformat{\subsection}{\large\bfseries\color{bleuDAE}}{\thesubsection}{0.6em}{}
\titleformat{\subsubsection}{\normalsize\bfseries\color{grisTexteDAE}}{\thesubsubsection}{0.6em}{}

\pagestyle{fancy}
\fancyhf{}
\fancyhead[L]{\small\textcolor{grisTexteDAE}{Projet GDAETF2 — Aide à la décision}}
\fancyhead[R]{\small\textcolor{grisTexteDAE}{Bibliothèque PDF}}
\fancyfoot[C]{\small\textcolor{grisTexteDAE}{\thepage\,/\,\pageref{LastPage}}}
\renewcommand{\headrulewidth}{0.3pt}
\renewcommand{\footrulewidth}{0pt}

\newcolumntype{Y}{>{\raggedright\arraybackslash}X}
\newcolumntype{Z}{>{\centering\arraybackslash}p{1.25cm}}
\newcommand{\recommande}{\textcolor{vertDAE}{\textbf{Recommandé}}}
\newcommand{\alternative}{\textcolor{orangeDAE}{\textbf{Alternative}}}
\newcommand{\note}[1]{\begin{tcolorbox}[enhanced,colback=bleuClairDAE,colframe=bleuDAE,boxrule=0.6pt,arc=1.5mm]#1\end{tcolorbox}}

\begin{document}

\begin{titlepage}
  \centering
  \vspace*{1.4cm}
  {\Large\textcolor{bleuDAE}{PROJET GDAETF2}\par}
  \vspace{1.1cm}
  {\Huge\bfseries\textcolor{bleuDAE}{Aide à la décision}\par}
  \vspace{0.35cm}
  {\LARGE\textcolor{grisTexteDAE}{Choix d'une bibliothèque PDF}\par}
  \vspace{1.2cm}
  \begin{tcolorbox}[enhanced,colback=bleuClairDAE,colframe=bleuDAE,boxrule=0.8pt,arc=2mm,width=0.88\textwidth]
    Comparaison des solutions utilisables avec CakePHP 5.3 pour préparer l'export PDF sécurisé des DAE prévu par l'ADR 0062.
  \end{tcolorbox}
  \vfill
  {\large\textcolor{vertDAE}{Choix recommandé : mPDF}\par}
  \vspace{0.3cm}
  {\small Document préparatoire — 6 octobre 2026\par}
\end{titlepage}

\tableofcontents
\clearpage

\section{Objet et contexte}

Ce document prépare le choix de la bibliothèque de génération PDF pour
l'application CakePHP 5.3. Le dépôt impose PHP \texttt{>= 8.2}, une génération
côté serveur, des dépendances hébergées dans l'environnement de l'application
et l'absence de conversion réalisée dans le navigateur.

L'\href{run:../../adr/0062-export-pdf-applicationform.md}{ADR 0062}
prévoit notamment :

\begin{itemize}
  \item une route protégée \texttt{/applicationforms/viewpdf/\{id\}} ;
  \item un service dédié et un gabarit de présentation local ;
  \item la réutilisation de la Policy et des \texttt{FieldAuthorizations} ;
  \item un tableau du cycle de validation, avec les statuts Accepté, Refusé et En attente ;
  \item un type MIME \texttt{application/pdf} et un nom de fichier explicite.
\end{itemize}

Le PDF ne doit donc pas constituer une seconde source de vérité pour les droits
ni nécessiter de modification du schéma de données.

\section{Critères de décision}

Les critères ci-dessous sont pondérés selon le besoin de l'application.

\begin{table}[htbp]
\centering
\small
\begin{tabularx}{0.9\textwidth}{>{\bfseries}p{0.34\textwidth}Z Y}
\toprule
\rowcolor{bleuDAE}
\textcolor{white}{Critère} & \textcolor{white}{Poids} & \textcolor{white}{Interprétation} \\
\midrule
Compatibilité PHP / CakePHP & 20\,\% & Installation et utilisation cohérentes avec PHP 8.2 et CakePHP 5.3. \\
Fidélité HTML / CSS & 20\,\% & Capacité à utiliser un gabarit local sans réécrire toute la mise en page. \\
Mise en page paginée & 15\,\% & Sauts de page, en-têtes, pieds de page et tableaux multipages. \\
Tables et workflow & 15\,\% & Rendu lisible du tableau des étapes et des votes. \\
Déploiement Docker & 15\,\% & Absence de binaire ou de runtime externe à maintenir. \\
Performance & 5\,\% & Temps de génération et consommation mémoire raisonnables. \\
Maintenance & 5\,\% & Pérennité du projet et clarté de son écosystème. \\
Testabilité & 5\,\% & Facilité de tester le contenu, la réponse HTTP et le rendu. \\
\bottomrule
\end{tabularx}
\caption{Critères et pondération.}
\label{tab:criteres}
\end{table}

La note de chaque solution est comprise entre 1 (défavorable) et 5
(excellent). Le score final est la moyenne pondérée des notes.

\section{Bibliothèques candidates}

\subsection{mPDF — \texttt{mpdf/mpdf}}

Bibliothèque PHP pure, conçue pour produire des documents paginés à partir
d'un sous-ensemble HTML/CSS. Elle convient particulièrement aux rapports, aux
tableaux, aux en-têtes et pieds de page, ainsi qu'aux documents comportant du
contenu français et des polices configurées localement.

\textbf{Forces :} pagination métier, tableaux, en-têtes, pieds de page,
intégration Composer et absence de dépendance système.

\textbf{Limites :} le support du CSS moderne n'est pas celui d'un navigateur
et la consommation mémoire doit être mesurée sur les DAE les plus volumineuses.

\subsection{Dompdf — \texttt{dompdf/dompdf}}

Bibliothèque PHP pure avec une intégration rapide pour un gabarit HTML/CSS
classique. Elle est adaptée à un document composé de blocs et de tableaux
simples.

\textbf{Forces :} simplicité, déploiement léger, API accessible et bon choix
pour un premier prototype.

\textbf{Limites :} support limité de Flexbox/Grid et contrôle moins prévisible
des mises en page complexes ou des sauts de page.

\subsection{TCPDF — \texttt{tecnickcom/tcpdf}}

Bibliothèque PHP pure robuste, particulièrement adaptée lorsque le document
est construit par programmation avec un contrôle précis des coordonnées et des
éléments PDF.

\textbf{Forces :} contrôle bas niveau, stabilité et faible dépendance à HTML.

\textbf{Limites :} le gabarit HTML prévu par l'ADR 0062 serait moins naturel à
exploiter et la mise en page demanderait davantage de code spécifique.

\subsection{wkhtmltopdf — \texttt{knplabs/knp-snappy}}

Solution fondée sur un binaire externe convertissant HTML/CSS via un moteur
WebKit. Le paquet PHP sert d'adaptateur, mais le binaire doit être installé et
maintenu dans l'image Docker.

\textbf{Forces :} très bon rendu du HTML/CSS classique et réutilisation de
pages proches de l'affichage navigateur.

\textbf{Limites :} moteur WebKit ancien, installation système, gestion des
processus et surface opérationnelle plus importantes.

\subsection{Browsershot — \texttt{spatie/browsershot}}

Solution pilotant Chromium via Node.js. Elle offre un rendu très proche d'un
navigateur moderne.

\textbf{Forces :} excellent support CSS moderne et rendu fidèle.

\textbf{Limites :} ajout de Node.js et Chromium, image Docker plus lourde,
temps de démarrage et exploitation plus complexes pour un document métier.

\subsection{WeasyPrint}

Moteur de composition PDF orienté documents imprimés, généralement utilisé via
Python et un processus externe depuis PHP.

\textbf{Forces :} modèle paginé et CSS d'impression de qualité.

\textbf{Limites :} absence d'intégration PHP native et ajout d'une chaîne
Python à l'exploitation CakePHP.

\subsection{FPDF — \texttt{setasign/fpdf}}

Bibliothèque légère de construction PDF programmatique.

\textbf{Forces :} simplicité, contrôle et faible empreinte.

\textbf{Limites :} aucune conversion HTML ; le positionnement, les retours à la
ligne et la pagination devraient être développés dans le service.

\begin{landscape}
\section{Matrice de décision}

\small
\begin{table}[htbp]
\centering
\setlength{\tabcolsep}{3.5pt}
\begin{tabularx}{\linewidth}{>{\raggedright\arraybackslash}p{0.22\linewidth}*{6}{>{\centering\arraybackslash}X}}
\toprule
\rowcolor{bleuDAE}
\textcolor{white}{Critère} & \textcolor{white}{mPDF} & \textcolor{white}{Dompdf} & \textcolor{white}{TCPDF} & \textcolor{white}{wkhtmltopdf} & \textcolor{white}{Browsershot} & \textcolor{white}{WeasyPrint} \\
\midrule
Compatibilité PHP / CakePHP & 5 & 5 & 5 & 4 & 3 & 2 \\
Fidélité HTML / CSS & 4 & 3 & 2 & 5 & 5 & 5 \\
Mise en page paginée & 5 & 3 & 4 & 4 & 4 & 5 \\
Tables et workflow & 5 & 4 & 3 & 4 & 4 & 5 \\
Déploiement Docker & 5 & 5 & 5 & 2 & 1 & 1 \\
Performance & 3 & 4 & 4 & 3 & 2 & 3 \\
Maintenance & 4 & 4 & 4 & 3 & 4 & 4 \\
Testabilité & 4 & 4 & 4 & 3 & 3 & 3 \\
\midrule
\rowcolor{vertClairDAE}
\textbf{Score pondéré / 5} & \textbf{4,55} & \textbf{3,95} & \textbf{3,65} & \textbf{3,65} & \textbf{3,25} & \textbf{3,65} \\
\bottomrule
\end{tabularx}
\caption{Matrice de décision pondérée.}
\label{tab:matrice}
\end{table}

\vspace{0.4cm}
\begin{tabularx}{\linewidth}{>{\bfseries}p{0.23\linewidth}Y}
\toprule
\rowcolor{bleuDAE}
\textcolor{white}{Solution} & \textcolor{white}{Lecture de la décision} \\
\midrule
mPDF & Meilleur compromis pour le besoin complet : documents paginés, tables, PHP pur et déploiement Docker simple. \\
Dompdf & Alternative raisonnable si le gabarit reste simple et si le prototype confirme la pagination. \\
TCPDF / FPDF & À réserver à un document entièrement programmatique, sans dépendre d'un gabarit HTML. \\
wkhtmltopdf / Browsershot / WeasyPrint & À réserver à un besoin fort de fidélité navigateur ou CSS d'impression justifiant des dépendances externes. \\
\bottomrule
\end{tabularx}
\end{landscape}

\section{Recommandation}

\note{\recommande{} de retenir \texttt{mpdf/mpdf} pour l'implantation, sous
réserve d'un prototype de rendu avec une DAE représentative. \texttt{dompdf/dompdf}
doit rester l'alternative si le gabarit retenu n'utilise que du HTML/CSS simple.}

mPDF respecte le mieux la combinaison des contraintes suivantes :

\begin{itemize}
  \item bibliothèque installable via Composer et exécutable en PHP pur ;
  \item absence de binaire ou de runtime externe dans le conteneur ;
  \item bonne prise en charge des tableaux et de la pagination ;
  \item séparation claire entre le service de collecte sécurisée et le gabarit ;
  \item tests possibles sans navigateur ni service distant.
\end{itemize}

Le choix ne dispense pas de vérifier les versions disponibles au moment de
l'ajout dans Composer, ni de mesurer la consommation mémoire sur les documents
les plus longs.

\section{Architecture d'intégration proposée}

La bibliothèque ne doit pas être appelée directement par le contrôleur au-delà
de l'orchestration HTTP. L'architecture cible est :

\begin{enumerate}
  \item \texttt{ApplicationformsController}\\
        \texttt{::viewpdf()} charge la DAE et applique
        \texttt{ApplicationformPolicy} avant tout rendu ;
  \item un service dédié collecte les données visibles pour l'identité courante,
        en réutilisant les autorisations existantes ;
  \item un gabarit local prépare le contenu HTML du document ;
  \item un adaptateur PDF encapsule mPDF et retourne le contenu binaire ;
  \item la réponse CakePHP définit \texttt{application/pdf}, le nom du fichier
        et l'ouverture dans un nouvel onglet côté interface.
\end{enumerate}

Aucune table, vue SQL ou migration n'est nécessaire pour produire le PDF.

\section{Plan de validation avant implantation}

\begin{enumerate}
  \item Ajouter un prototype avec une DAE comportant toutes les zones visibles.
  \item Vérifier les champs autorisés et l'absence réelle des champs interdits.
  \item Vérifier les trois statuts du workflow et les commentaires autorisés.
  \item Vérifier les sauts de page, l'en-tête, le pied de page et le tableau multipage.
  \item Vérifier les accents, les polices, les montants et les dates en français.
  \item Mesurer le temps de génération et la mémoire sur une DAE volumineuse.
  \item Ajouter les tests unitaires du service et les tests HTTP de la réponse PDF.
  \item Documenter le choix définitif dans l'ADR approprié avant d'ajouter la dépendance Composer.
\end{enumerate}

\section{Conclusion}

Pour la première version, mPDF offre le meilleur équilibre entre qualité de
document métier, intégration PHP, déploiement et testabilité. Dompdf constitue
une alternative valable si le prototype montre que les exigences de pagination
restent simples. Les moteurs fondés sur un navigateur ou un processus externe
ne sont justifiés que par un besoin de rendu qui ne pourrait pas être couvert
par une bibliothèque PHP pure.

\end{document}
